% GENERATED FILE. Edit canonical lessons, then run npm run generate:books.
% canonical-source-hash: fnv1a64:f3a6d5615671ffad
% canonical-lessons: HI-C140-angur, HI-C140-santara, HI-C140-tarbuz, HI-C140-nariyal, HI-C140-dahi

\chapter{Grapes, Orange, Watermelon, Coconut, Yoghurt}
\label{ch:hi-grapes-orange-watermelon-coconut-yoghurt}
\hlchaptermodality{140}

\begin{chapteropening}
\textbf{By the end of this chapter:} \emph{I can say grapes, an orange, a watermelon, a coconut and yoghurt.}

Complete the last lesson of chapter 140: i can say grapes, an orange, a watermelon, a coconut and yoghurt.
\end{chapteropening}

\section[aṅgūr]{\hi{अंगूर} (aṅgūr) — grapes}

\label{lesson:HI-C140-angur}



\begin{quote}
Before the new one: say the Hindi for a banana, then the Hindi for an apple.
\end{quote}

\subsection*{You'll want to know: \hi{अंगूर}}
\textbf{\hi{अंगूर}} — \emph{aṅgūr} — \textquotedblleft{}grapes\textquotedblright{}.

\begin{grammarlens}[title={What you've built}]
One more word for this chapter.
\end{grammarlens}

\subsection*{Your turn}
\begin{itemize}
\raggedright
  \item \emph{Say it:} \emph{aṅgūr}
  \item \emph{Say it:} \emph{aṅgūr}, once more
  \item \emph{Say it:} say \emph{aṅgūr}
  \item \emph{Recall:} say the Hindi for a banana, then the Hindi for an apple, then say \emph{aṅgūr} again
\end{itemize}

\begin{tcolorbox}[breakable,colback=teal!4,colframe=teal!35!black,title={Before you move on}]
What does \hi{अंगूर} mean? (\textquotedblleft{}Grapes\textquotedblright{}.) Say it once more.
\end{tcolorbox}
\section[santarā]{\hi{संतरा} (santarā) — an orange}

\label{lesson:HI-C140-santara}



\begin{quote}
Before the new one: say the Hindi for an apple, then the Hindi for grapes.

Count \textbf{\hi{छह}} to \textbf{\hi{दस}}, then \textbf{\hi{ग्यारह}} to \textbf{\hi{बीस}}, and by tens \textbf{\hi{सत्तर}, \hi{अस्सी}, \hi{नब्बे}}. Which one is \emph{one less than twenty}? (\textbf{\hi{उन्नीस}}.)
\end{quote}

\subsection*{You'll want to know: \hi{संतरा}}
\textbf{\hi{संतरा}} — \emph{santarā} — \textquotedblleft{}an orange\textquotedblright{}.

\begin{grammarlens}[title={What you've built}]
One more word for this chapter.
\end{grammarlens}

\subsection*{Your turn}
\begin{itemize}
\raggedright
  \item \emph{Say it:} \emph{santarā}
  \item \emph{Say it:} \emph{santarā}, once more
  \item \emph{Say it:} say \emph{santarā}
  \item \emph{Recall:} say the Hindi for an apple, then the Hindi for grapes, then say \emph{santarā} again
\end{itemize}

\begin{tcolorbox}[breakable,colback=teal!4,colframe=teal!35!black,title={Before you move on}]
What does \hi{संतरा} mean? (\textquotedblleft{}An orange\textquotedblright{}.) Say it once more.
\end{tcolorbox}
\section[tarbūz]{\hi{तरबूज़} (tarbūz) — a watermelon}

\label{lesson:HI-C140-tarbuz}



\begin{quote}
Before the new one: say the Hindi for grapes, then the Hindi for an orange.
\end{quote}

\subsection*{You'll want to know: \hi{तरबूज़}}
\textbf{\hi{तरबूज़}} — \emph{tarbūz} — \textquotedblleft{}a watermelon\textquotedblright{}.

\begin{grammarlens}[title={What you've built}]
One more word for this chapter.
\end{grammarlens}

\subsection*{Your turn}
\begin{itemize}
\raggedright
  \item \emph{Say it:} \emph{tarbūz}
  \item \emph{Say it:} \emph{tarbūz}, once more
  \item \emph{Say it:} say \emph{tarbūz}
  \item \emph{Recall:} say the Hindi for grapes, then the Hindi for an orange, then say \emph{tarbūz} again
\end{itemize}

\begin{tcolorbox}[breakable,colback=teal!4,colframe=teal!35!black,title={Before you move on}]
What does \hi{तरबूज़} mean? (\textquotedblleft{}A watermelon\textquotedblright{}.) Say it once more.
\end{tcolorbox}
\section[nāriyal]{\hi{नारियल} (nāriyal) — a coconut}

\label{lesson:HI-C140-nariyal}



\begin{quote}
Before the new one: say the Hindi for an orange, then the Hindi for a watermelon.
\end{quote}

\subsection*{You'll want to know: \hi{नारियल}}
\textbf{\hi{नारियल}} — \emph{nāriyal} — \textquotedblleft{}a coconut\textquotedblright{}.

\begin{grammarlens}[title={What you've built}]
One more word for this chapter.
\end{grammarlens}

\subsection*{Your turn}
\begin{itemize}
\raggedright
  \item \emph{Say it:} \emph{nāriyal}
  \item \emph{Say it:} \emph{nāriyal}, once more
  \item \emph{Say it:} say \emph{nāriyal}
  \item \emph{Recall:} say the Hindi for an orange, then the Hindi for a watermelon, then say \emph{nāriyal} again
\end{itemize}

\begin{tcolorbox}[breakable,colback=teal!4,colframe=teal!35!black,title={Before you move on}]
What does \hi{नारियल} mean? (\textquotedblleft{}A coconut\textquotedblright{}.) Say it once more.
\end{tcolorbox}
\section[dahī]{\hi{दही} (dahī) — yoghurt}

\label{lesson:HI-C140-dahi}



\begin{quote}
Before the new one: say the Hindi for a watermelon, then the Hindi for a coconut.
\end{quote}

\subsection*{You'll want to know: \hi{दही}}
\textbf{\hi{दही}} — \emph{dahī} — \textquotedblleft{}yoghurt\textquotedblright{}.

\begin{grammarlens}[title={What you've built}]
One more word for this chapter.
\end{grammarlens}

\subsection*{Your turn}
\begin{itemize}
\raggedright
  \item \emph{Say it:} \emph{dahī}
  \item \emph{Say it:} \emph{dahī}, once more
  \item \emph{Say it:} say \emph{dahī}
  \item \emph{Recall:} say the Hindi for a watermelon, then the Hindi for a coconut, then say \emph{dahī} again
\end{itemize}

\begin{tcolorbox}[breakable,colback=teal!4,colframe=teal!35!black,title={Before you move on}]
What does \hi{दही} mean? (\textquotedblleft{}Yoghurt\textquotedblright{}.) Say it once more.
\end{tcolorbox}
